\subsection{The full $E_6$ involution ladder closes on the cubic $36$--$45$ incidence geometry}
\label{sec:20260917-e6-involution-ladder-flags}

The preceding involution reconstruction can be sharpened so that all four
nontrivial involution degrees of $W(E_6)$ acquire direct cubic-surface meaning.
The degree-one class is the $36$ reflections, already identified objectwise with
the $36$ double-sixes through their moved twelve-line supports.

For two distinct double-sixes $D_1,D_2$, their supports meet in either four or
six cubic lines, with exact census
\[
 \binom{36}{2}=270+360.
\]
The corresponding reflections commute if and only if the overlap is four.
Their $270$ products are distinct and exhaust the degree-two involution class.
Thus
\[
 \boxed{\{d=2\text{ involutions}\}
 \longleftrightarrow
 \{\text{syzygetic double-six pairs}\}},
 \qquad |D_1\cap D_2|=4.
\]
The remaining $360$ pairs have six-line overlap and are noncommuting.

At the opposite end, each degree-four involution fixes exactly three cubic
lines pointwise, those three lines form a tritangent, and the resulting map is
a bijection
\[
 \boxed{\{d=4\text{ involutions}\}\longleftrightarrow
        \{45\text{ tritangents}\}}.
\]
Write $q_T$ for the unique degree-four involution attached to tritangent $T$.

Now let $w$ be one of the $540$ degree-three involutions.  Its three fixed
lines form a tritangent $T$, and $w$ commutes with $q_T$.  The product
\[
 r=wq_T
\]
is a degree-one reflection.  If $D$ is the double-six moved by $r$, then
\[
 T\cap D=\varnothing.
\]
For each fixed $T$, the twelve degree-three involutions above $T$ map to the
twelve double-sixes disjoint from $T$, each exactly once.  Globally this gives
the canonical bijection
\[
 \boxed{
 \{d=3\text{ involutions}\}
 \longleftrightarrow
 \{(T,D):T\text{ tritangent},\ D\text{ double-six},\ T\cap D=\varnothing\}},
\]
and therefore
\[
 \boxed{540=45\cdot12}.
\]
So the four involution degrees organize the cubic geometry as
\[
 36\ \longrightarrow\ 270\ \longrightarrow\ 540\ \longrightarrow\ 45,
\]
with carriers double-sixes, syzygetic double-six pairs, disjoint
tritangent--double-six flags, and tritangents respectively.  The complementary
$1080$ tritangent/double-six pairs have two-line intersection, giving again the
Pass~4659 census
\[
 45\cdot36=540+1080.
\]

The executable certificate is
\texttt{analysis/w33\_e6\_involution\_ladder\_flags.py}, with frozen output
\texttt{data/w33\_e6\_involution\_ladder\_flags.json}.  This is an exact finite
$W(E_6)$/cubic-surface theorem.  Transport to the W33/Suzuki carriers uses the
separately certified equivariant bridges and does not by itself assert a
physical $E_6$ gauge symmetry.
